\documentclass[11pt]{article}
\usepackage[a4paper,margin=1in]{geometry}
\usepackage[T1]{fontenc}
\usepackage[utf8]{inputenc}
\usepackage{lmodern}
\usepackage{microtype}
\usepackage{amsmath,amssymb,amsthm,mathtools}
\usepackage{enumitem}
\usepackage[hidelinks]{hyperref}
\usepackage{titlesec}
\titleformat{\section}{\large\bfseries}{\thesection}{0.6em}{}
\titleformat{\subsection}{\normalsize\bfseries}{\thesubsection}{0.6em}{}
\newtheorem*{satz}{Theorem}
\newtheorem*{bemerkung}{Remark}
\DeclareMathOperator{\Fix}{Fix}
\DeclareMathOperator{\Gal}{Gal}
\DeclareMathOperator{\disc}{disc}
\newcommand{\webersection}[2]{\section*{\S #1. #2}\addcontentsline{toc}{section}{\S #1. #2}}
\setlength{\parindent}{0pt}
\setlength{\parskip}{0.6em}

\title{Heinrich Weber\\\emph{Lehrbuch der Algebra}\\Working English Translation part 09\\Binary Linear Substitutions and Invariants}
\author{Translated into modern mathematical English}
\date{}
\begin{document}
\maketitle

\section*{Translator's note}
This working translation continues the cleaned Weber source after the material on permutation groups as linear substitution groups.  The notation follows the source closely, while spelling and terminology are normalized to current mathematical English.

\webersection{40}{Groups of binary linear substitutions with real coefficients}

For the binary substitutions
\begin{equation}
 y_1=\alpha x_1+\beta x_2,\qquad
 y_2=\gamma x_1+\delta x_2,
\end{equation}
the determinant is
\[
D=\alpha\delta-\beta\gamma .
\]
We shall first suppose that the coefficients
\(\alpha,\beta,\gamma,\delta\) are real, and ask under what conditions a group of such substitutions can be finite.

The characteristic equation of the substitution is
\begin{equation}
\begin{vmatrix}
\alpha-\lambda & \beta\\
\gamma & \delta-\lambda
\end{vmatrix}
=
\lambda^2-(\alpha+\delta)\lambda+D=0.
\end{equation}
Its roots are
\begin{equation}
\lambda=
\frac{\alpha+\delta\pm\sqrt{(\alpha+\delta)^2-4D}}{2}.
\end{equation}

If
\[
(\alpha+\delta)^2-4D>0,
\]
then the roots are real and distinct.  In that case the substitution can be brought, by a similar substitution, into the form
\begin{equation}
y_1=\lambda_1x_1,\qquad y_2=\lambda_2x_2.
\end{equation}
The \(n\)-fold iterate is then
\[
y_1=\lambda_1^n x_1,\qquad y_2=\lambda_2^n x_2.
\]
For some positive power to become the identity, \(\lambda_1\) and \(\lambda_2\) must be roots of unity.  Since they are real, this forces
\[
\lambda_1,\lambda_2\in\{1,-1\}.
\]

If
\[
(\alpha+\delta)^2-4D=0,
\]
then the two characteristic roots coincide: \(\lambda_1=\lambda_2=\lambda\).  The substitution can then either be reduced to the diagonal form above, or to the Jordan form
\begin{equation}
y_1=\lambda x_1,\qquad y_2=x_1+\lambda x_2.
\end{equation}
In the latter case the \(n\)-fold iterate contains the term
\[
y_2=n\lambda^{n-1}x_1+\lambda^n x_2,
\]
and therefore cannot be the identity for any positive \(n\), unless the nilpotent part vanishes.  Thus a finite group contains no nontrivial substitution of this Jordan type.

It remains to consider the case
\[
(\alpha+\delta)^2-4D<0.
\]
Then the two characteristic roots are conjugate complex roots.  A real substitution with such roots is, over the real field, similar to a rotation-dilatation.  If the substitution has finite order, its characteristic roots must again be roots of unity; hence the quotient of the two roots is a root of unity and the corresponding projective transformation is periodic.

Consequently, a real binary linear substitution which belongs to a finite group is characterized by the fact that, after multiplication by a scalar if necessary, it can be interpreted as a rotation through an angle commensurable with \(2\pi\), or as one of the exceptional involutory cases in which the characteristic roots are \(1\) and \(-1\).

This gives a first reduction of the problem of finite real binary substitution groups to the problem of finite rotation groups and their associated projective actions.

\webersection{41}{Invariants and covariants of groups of linear substitutions}

The theory of linear substitutions is connected in the closest way with the theory of algebraic forms.  By a \emph{form} in the variables
\[
x_1,x_2,\ldots,x_n
\]
we mean an integral homogeneous function of these variables.  A form
\[
f(x_1,\ldots,x_n)
\]
of degree \(p\) is transformed by a linear substitution \(A\) into another form
\[
f'(x'_1,\ldots,x'_n)
\]
of the same degree in the new variables.  If \(f'=f\), that is, if the form is transformed into itself by the substitution, then \(f\) is called an \emph{invariant} of \(A\).

More generally, suppose that \(f\) is a function not only of the variables \(x\), but also of other variables \(u,v,\ldots\).  If \(f\) is transformed into itself by applying \(A\) to the \(x\)-variables and at the same time applying corresponding substitutions to the variables \(u,v,\ldots\), then \(f\) is called a \emph{covariant}.

We now consider a group \(S\) of linear substitutions in the variables
\[
x_1,\ldots,x_n,
\]
and seek those forms which are simultaneously invariants, or covariants, for all substitutions of the group.

Let
\[
F_1,F_2,\ldots,F_r
\]
be a system of such invariants which are rationally independent, meaning that no identical rational relation holds among them, while every other invariant can be rationally expressed through them.  Such a system is called a \emph{basis} of the invariant system of the group \(S\).

The most general invariant has the form
\begin{equation}
F=A_1F_1+A_2F_2+\cdots+A_rF_r,
\end{equation}
where \(A_1,A_2,\ldots,A_r\) are arbitrary constants, or more generally quantities lying in the appropriate coefficient field.

From the general theory of algebraic forms one obtains the following fundamental facts.

\begin{enumerate}[label=\Roman*.]
\item There always exists a finite system of invariants from which all the remaining invariants may be generated rationally.
\item If the group is finite, the invariant system may be chosen so that it reflects the finite averaging process over the substitutions of the group.
\item The corresponding covariants are obtained in the same manner once the simultaneous transformation law of the auxiliary variables is fixed.
\end{enumerate}

These facts allow the invariant theory of a finite linear substitution group to be treated algebraically: the full system of invariant forms is replaced by a finite basis and the rational relations among its elements.

\end{document}
